% Drawing 138/01-2 -- folder 138, batch 01, archive page 19 (PDF page 20).
% Transcription: à gauche, un arc isolé ; plus bas, une ovale traversée par des arcs joignant $S_0$ (en haut) à $S_\infty$ (en bas), passant par des points marqués $R'$, $S_1$, $E$, $E'$, $R''$ sur l'ovale
% Kind: oval crossed by arcs (sphere with paths). Method 2 (schematic redraw from the facsimile). Status: partial.
% Notes: Left: an isolated arc. Below: an oval crossed by arcs from S_0 (top) to S_infinity (bottom). The marked points on the oval are crowded: R' on the upper edge, then S_1, E, E' on the lower edge, R'' on an arc below; drawn as three arcs through them. How many arcs he drew (three or four) and which arc passes through which point are a reading. The T and the text at the right belong to the transcription.
% Coordinates: 1 unit = 1 mm = 1/100 of the crop width (bbox in page fractions: [0.14, 0.27, 0.34, 0.53]).
% Redrawn by Claude (claude-opus-5-5), issue #73, 2026-11-07; not checked by a human; not compiled (no TeX in the session).
% Rights: reproduces part of the work -- behind RENDER_TRANSCRIPTS, never published by itself.
\documentclass[tikz,border=3pt]{standalone}
\usepackage{amsmath,amssymb}
\usetikzlibrary{patterns,arrows.meta}
\tikzset{pencil/.style={black!45}, struck/.style={black!35}, mark/.style={circle,draw,inner sep=1.3pt}, pt/.style={circle,fill,inner sep=0.9pt}, opt/.style={circle,draw,fill=white,inner sep=1.1pt}}
\begin{document}
\begin{tikzpicture}[x=1mm,y=1mm,line cap=round,line join=round,font=\scriptsize]
\draw[very thick] plot[smooth, tension=0.55] coordinates {(6,-43) (7,-26) (12,-14) (20,-7) (28,-2)};
\draw plot[smooth, tension=0.55] coordinates {(10,-110) (30,-103) (55,-100) (70,-104) (80,-113)};
\draw plot[smooth, tension=0.55] coordinates {(10,-110) (4,-118) (6,-126) (25,-131) (47,-131) (65,-126) (82,-120)};
\draw plot[smooth, tension=0.55] coordinates {(41,-83) (30,-92) (25,-101) (24,-118) (25,-129) (28,-145) (31,-152) (40,-165) (51,-173)};
\draw plot[smooth, tension=0.55] coordinates {(41,-83) (37,-100) (35,-128) (38,-150) (51,-173)};
\draw plot[smooth, tension=0.55] coordinates {(41,-83) (42,-105) (45,-129) (47,-150) (51,-173)};
\node[pt] at (41,-83) {};
\node[pt] at (25,-101) {};
\node[pt] at (25,-129) {};
\node[pt] at (35,-129) {};
\node[pt] at (47,-130) {};
\node[pt] at (31,-152) {};
\node[pt] at (51,-173) {};
\node[inner sep=1pt] at (47,-75) {\scriptsize $S_0$};
\node[inner sep=1pt] at (20,-96) {\scriptsize $R'$};
\node[inner sep=1pt] at (18,-128) {\tiny $S_1$};
\node[inner sep=1pt] at (32,-140) {\tiny $E$};
\node[inner sep=1pt] at (43,-139) {\tiny $E'$};
\node[inner sep=1pt] at (26,-158) {\tiny $R''$};
\node[inner sep=1pt] at (56,-182) {\scriptsize $S_\infty$};
\end{tikzpicture}
\end{document}
